%% ch05 --- Jedna linijka algebry, każde możliwe zachowanie
%%
%% Zalążek, wygenerowany przez tools/gen_stubs.py z tools/chapters.json.
%% Poniższy opis jest kontraktem tego rozdziału: co ma uzasadnić, co ma dać
%% czytelnikowi i czego NIE ma zabierać sąsiednim rozdziałom. Napisanie
%% rozdziału oznacza usunięcie zalążkowego bloku (tego, który zaczyna się po
%% linii \chapter) i zastąpienie go rozdziałem -- po czym ten plik nie jest
%% już generowany.

\chapter{Jedna linijka algebry, każde możliwe zachowanie}
\label{ch:logistic}

\chapterstub{%
\textit{[Opis do przetłumaczenia --- patrz wydanie angielskie.]}

Argument: the discrete logistic rule \code{x[n+1] = r*x[n]*(1-x[n])} \dash{}
one multiplication short of the linear rules of Chapter 2 \dash{} produces
every behaviour the book is about as the single knob r is turned. r = 2.8
settles to a fixed point; r = 3.2 alternates between two values for ever; r
= 3.5 cycles through four; r = 3.9 never repeats and never settles, although
every value is produced by the same exact rule. The cobweb diagram from
Chapter 2, now on a parabola, shows why: the fixed point's slope passes
minus one and the fixed point loses its grip. Robert May's 1976 paper in
Nature, written for ecologists, is the historical anchor, and its plea that
everyone should meet this map early is the chapter's reason to exist.
Payoff: the reader can iterate the map, draw its cobweb, find the fixed
point \code{1 - 1/r} and say for which r it is stable \dash{} and has seen
irregular, never-repeating output come from a rule a child can compute.
Traps to elicit: a simple rule gives simple behaviour; irregular output
needs a random or complicated cause; if it has not repeated after a thousand
steps it is about to. Listings: iterating the map at four values of r, in
pure Python so every digit is reproducible; the cobweb data; period
detection. Labs: detect the period of an orbit; find r where the fixed point
loses stability. Figures: four time series; four cobwebs.

\medskip
\textbf{Budżet słów:} 3000.
}
